%! TEX root = main.tex

\section{Conservation of Energy}
\label{sc:22}

The conservation of mechanical energy under periodic boundary conditions is established by showing that the time derivative of the prescribed energy functional vanishes identically.

\subsection{Definition of the Energy Functional}

The energy functional of the shallow water system over the spatial domain \((0,L)\) is defined as:

\begin{equation}
E(t) = \int_{0}^{L} \left( \frac{1}{2}g\eta^{2} + \frac{1}{2}\frac{q^{2}}{H} \right) dx
\label{eq:energy}
\end{equation}

where the first term represents the potential energy associated with the free surface elevation \(\eta\), and the second term represents the kinetic energy associated with the discharge \(q\). To prove conservation, it suffices to show that \(\frac{dE}{dt} = 0\).

\subsection{Time Derivative of the Energy}

Assuming sufficient smoothness of \(\eta(x,t)\) and \(q(x,t)\), the Leibniz rule permits the exchange of differentiation and integration:

\begin{equation}
\frac{dE}{dt} = \frac{d}{dt} \int_{0}^{L} \left( \frac{1}{2}g\eta^{2} + \frac{1}{2}\frac{q^{2}}{H} \right) dx = \int_{0}^{L} \frac{\partial}{\partial t} \left( \frac{1}{2}g\eta^{2} + \frac{1}{2}\frac{q^{2}}{H} \right) dx
\end{equation}

Applying the chain rule to compute the partial time derivatives of the integrand yields:

\begin{equation}
\frac{dE}{dt} = \int_{0}^{L} \left( g \eta \eta_t + \frac{1}{H} q q_t \right) dx
\label{eq:energy_dt}
\end{equation}

\subsection{Substitution of the Strong Formulation}

Recalling the linearized shallow water equations \eqref{eq:strong_system}:

\begin{equation}
\begin{cases}
\eta_{t} = -q_{x} \\
q_{t} = -gH\eta_{x}
\end{cases}
\label{eq:strong_substituted}
\end{equation}

Substituting these expressions for \(\eta_t\) and \(q_t\) into \eqref{eq:energy_dt}:

\begin{equation}
\frac{dE}{dt} = \int_{0}^{L} \left( g \eta (-q_x) + \frac{1}{H} q (-gH\eta_x) \right) dx
\end{equation}

\subsection{Algebraic Simplification and the Product Rule}

Factoring out the common constant \(-g\), the integral simplifies to:

\begin{equation}
\frac{dE}{dt} = -g \int_{0}^{L} \left( \eta q_x + q \eta_x \right) dx
\end{equation}

The integrand is recognized as the exact spatial derivative of the product \(\eta q\) via the product rule \(\frac{\partial}{\partial x}(\eta q) = \eta_x q + \eta q_x\). Therefore:

\begin{equation}
\frac{dE}{dt} = -g \int_{0}^{L} \frac{\partial}{\partial x} (\eta q) dx
\end{equation}

\subsection{Application of the Fundamental Theorem of Calculus}

Integrating the exact derivative over \([0,L]\) yields the evaluation of the boundary terms:

\begin{equation}
\frac{dE}{dt} = -g \Big[ \eta(x,t) q(x,t) \Big]_{x=0}^{x=L} = -g \left( \eta(L,t)q(L,t) - \eta(0,t)q(0,t) \right)
\end{equation}

\subsection{Application of Boundary Conditions}

The periodic boundary conditions prescribe \(\eta(0,t) = \eta(L,t)\) and \(q(0,t) = q(L,t)\). Substituting these into the boundary evaluation:

\begin{equation}
\eta(L,t)q(L,t) - \eta(0,t)q(0,t) = \eta(L,t)q(L,t) - \eta(L,t)q(L,t) = 0
\end{equation}

\subsection{Energy Conservation}

Therefore, the time derivative of the energy is exactly zero:

\begin{equation}
\frac{dE}{dt} = 0
\end{equation}

This rigorously demonstrates that the energy \(E(t)\) defined in \eqref{eq:energy} is conserved for all \(t \in (0,T)\) under the prescribed periodic boundary conditions.

% REMARK : USEFUL FOR ORAL EXAMINATION %
% Note that this result was established using the first-order system \((\eta, q)\) rather than the second-order wave equation from section \ref{sc:21}. The energy functional \eqref{eq:energy} explicitly tracks both the potential energy (proportional to \(\eta^2\)) and the kinetic energy (proportional to \(q^2/H\)). The use of the coupled first-order system directly links the strong-form PDE operators to the energy structure.